\ifajar
  \section{Latihan}
\else
  \section{Bab 13: Sistem Real-Time dan Asinkron}
\fi

Seluruh file yang dipakai pada latihan ini tersedia di
\berkas{sources/chapter13/}, dihitung relatif terhadap akar repositori.
Latihan ini hanya membutuhkan Python 3, tanpa Docker dan tanpa pustaka
tambahan, termasuk WebSocket-nya.

Seluruh permintaan diarahkan ke layanan yang berjalan di komputer sendiri,
dan sifatnya mengamati, bukan membebani. Angka milidetik tiap mahasiswa akan
berbeda, sedangkan pola perbandingannya dan jumlah percobaannya akan sama.

Penyiapannya tiga perintah, dan urutannya berpengaruh, seperti pada Listing
\ref{lst:c13-siapkan}.

\begin{lstlisting}[language=bash, caption={Menyalakan ketiga layanan, tersedia di \berkas{sources/chapter13/}}, label={lst:c13-siapkan}]
cd sources/chapter13
python3 aplikasi.py &
python3 webhook.py &
python3 websocket.py &

# Keluarannya:
# Melayani di http://127.0.0.1:8100/sse dan /polling  (Ctrl+C untuk berhenti)
# Webhook di http://127.0.0.1:8102/webhook  (Ctrl+C untuk berhenti)
# WebSocket di ws://127.0.0.1:8101/  (Ctrl+C untuk berhenti)
\end{lstlisting}

\ifajar
  \subsection*{Latihan 1: Membandingkan Tiga Transport}
\else
  \subsection*{Prosedur 1: Membandingkan Tiga Transport}
\fi

Latihan ini mengukur keterlambatan dan biaya ketiga transport pada peristiwa
yang sama, lalu memisahkan biaya yang terlihat dari biaya yang tersembunyi.

\begin{enumerate}
  \item Jalankan \berkas{ukur-transport.py}, seperti pada Listing
        \ref{lst:c13-lat1-jalan}, lalu catat ketiga barisnya.
  \item Catat jumlah permintaan tiap transport, lalu hitung berapa permintaan
        yang \textit{response}-nya kosong pada varian \textit{polling}.
  \item Buka \texttt{http://127.0.0.1:8100/sse} dengan \texttt{curl -sN},
        lalu terbitkan satu peristiwa dari terminal lain. Catat bentuk blok
        yang tiba.
  \item Isi Tabel \ref{tab:c13-lat1-hasil}. Tabel \ref{tab:c13-lat1-contoh}
        memperlihatkan contoh yang sudah terisi.
  \item Jawab tiga pertanyaan berikut. Pertama, mengapa keterlambatan
        \textit{polling} tidak membaik walaupun servernya cepat? Kedua,
        mengapa kolom byte tidak cukup untuk menilai biaya \textit{polling}?
        Ketiga, kapan WebSocket pantas dipilih walaupun SSE lebih sederhana?
\end{enumerate}

\begin{lstlisting}[language=bash, caption={Mengukur ketiga transport, tersedia di \berkas{sources/chapter13/ukur-transport.py}}, label={lst:c13-lat1-jalan}]
python3 ukur-transport.py

# Keluarannya:
#   Transport   Peristiwa  Tengah ms  Tertinggi     Byte  Permintaan
#   polling            10      199.8      202.0      639          11
#   sse                10        1.8        2.0      400           1
#   websocket          10        1.9        2.1      490           1
\end{lstlisting}

\begin{table}[!htb]
  \centering
  \caption{Contoh hasil Latihan 1 yang sudah terisi}
  \label{tab:c13-lat1-contoh}
  \begin{tabularx}{\textwidth}{Xlll}
    \hline
    \textbf{Transport} & \textbf{Tengah ms} & \textbf{Byte} & \textbf{Permintaan} \\
    \hline
    \textit{polling} 0,5 detik & 199,8 & 639 & 11 \\
    \hline
    SSE & 1,8 & 400 & 1 \\
    \hline
    WebSocket & 1,9 & 490 & 1 \\
    \hline
  \end{tabularx}
\end{table}

\begin{table}[!htb]
  \centering
  \caption{Lembar isian hasil pengukuran Latihan 1}
  \label{tab:c13-lat1-hasil}
  \begin{tabularx}{\textwidth}{Xlll}
    \hline
    \textbf{Transport} & \textbf{Tengah ms} & \textbf{Byte} & \textbf{Permintaan} \\
    \hline
    \textit{polling} 0,5 detik & \dots & \dots & \dots \\
    \hline
    SSE & \dots & \dots & \dots \\
    \hline
    WebSocket & \dots & \dots & \dots \\
    \hline
  \end{tabularx}
\end{table}

\ifajar
  \subsection*{Latihan 2: Jeda Polling dan Keterlambatannya}
\else
  \subsection*{Prosedur 2: Jeda Polling dan Keterlambatannya}
\fi

Latihan ini membuktikan bahwa keterlambatan \textit{polling} ditentukan
jedanya, dan bahwa memperkecil jeda membayar dengan jumlah permintaan.

\begin{enumerate}
  \item Ubah \texttt{JEDA\_POLLING\_DETIK} di \berkas{ukur-transport.py}
        menjadi 0,1 lalu jalankan pengukurannya. Catat keterlambatan dan
        jumlah permintaannya.
  \item Ulangi dengan nilai 0,5 dan 1,0.
  \item Bandingkan ketiganya dengan selisih jeda penerbitan peristiwa, yaitu
        \texttt{JEDA\_ANTARPENERBITAN\_DETIK} yang bernilai 0,3 detik.
  \item Isi Tabel \ref{tab:c13-lat2-hasil}. Tabel \ref{tab:c13-lat2-contoh}
        memperlihatkan contoh yang sudah terisi.
  \item Jawab dua pertanyaan berikut. Pertama, mengapa keterlambatannya
        mendekati jeda \textit{polling} dikurangi jeda penerbitan, bukan
        setengah jeda \textit{polling}? Kedua, pada beban nyata dengan
        peristiwa yang datang sembarang waktu, berapa keterlambatan rata-rata
        yang pantas diharapkan?
\end{enumerate}

\begin{lstlisting}[language=bash, caption={Mengukur satu transport pada beberapa jeda, memakai \berkas{sources/chapter13/ukur-transport.py}}, label={lst:c13-lat2-jalan}]
# JEDA_POLLING_DETIK diubah, lalu pengukurannya dijalankan ulang:
python3 ukur-transport.py

# Baris polling pada tiga nilai jeda, 0.1 lalu 0.5 lalu 1.0 detik:
#   polling            10        3.3        5.3      997          31
#   polling            10      200.2      202.6      657          11
#   polling            10      700.1      701.0      657          11
\end{lstlisting}

\begin{table}[!htb]
  \centering
  \caption{Contoh hasil Latihan 2 yang sudah terisi}
  \label{tab:c13-lat2-contoh}
  \begin{tabularx}{\textwidth}{Xlll}
    \hline
    \textbf{Jeda polling} & \textbf{Tengah ms} & \textbf{Byte} & \textbf{Permintaan} \\
    \hline
    0,1 detik & 3,3 & 997 & 31 \\
    \hline
    0,5 detik & 200,2 & 657 & 11 \\
    \hline
    1,0 detik & 700,1 & 657 & 11 \\
    \hline
  \end{tabularx}
\end{table}

\begin{table}[!htb]
  \centering
  \caption{Lembar isian hasil pengukuran Latihan 2}
  \label{tab:c13-lat2-hasil}
  \begin{tabularx}{\textwidth}{Xlll}
    \hline
    \textbf{Jeda polling} & \textbf{Tengah ms} & \textbf{Byte} & \textbf{Permintaan} \\
    \hline
    0,1 detik & \dots & \dots & \dots \\
    \hline
    0,5 detik & \dots & \dots & \dots \\
    \hline
    1,0 detik & \dots & \dots & \dots \\
    \hline
  \end{tabularx}
\end{table}

\ifajar
  \subsection*{Latihan 3: Retry dan Jeda yang Berlipat}
\else
  \subsection*{Prosedur 3: Retry dan Jeda yang Berlipat}
\fi

Latihan ini memperlihatkan jeda \textit{retry} bekerja, dan membuktikan
angkanya dapat dihitung sebelum diukur.

\begin{enumerate}
  \item Jalankan \berkas{ukur-antrean.py}, seperti pada Listing
        \ref{lst:c13-lat3-jalan}, lalu catat keenam percobaannya beserta
        detiknya.
  \item Hitung selisih detik antarpercobaan \texttt{kerja-0002}, lalu
        bandingkan dengan \texttt{JEDA\_DASAR\_DETIK} di
        \berkas{antrean.py}.
  \item Ubah \texttt{JEDA\_DASAR\_DETIK} menjadi 1,0 lalu ulangi, dan catat
        detik percobaan kedua dan ketiganya.
  \item Isi Tabel \ref{tab:c13-lat3-hasil}. Tabel \ref{tab:c13-lat3-contoh}
        memperlihatkan contoh yang sudah terisi.
  \item Jawab tiga pertanyaan berikut. Pertama, berapa lama satu pekerjaan
        bertahan sebelum percobaan kelimanya, bila jeda dasarnya 0,5 detik?
        Kedua, mengapa \textit{retry} tanpa jeda memperburuk keadaan? Ketiga,
        ke mana pekerjaan yang tetap gagal sesudah batas percobaan pantas
        dipindahkan?
\end{enumerate}

\begin{lstlisting}[language=bash, caption={Menjalankan satu daur antrean, tersedia di \berkas{sources/chapter13/ukur-antrean.py}}, label={lst:c13-lat3-jalan}]
python3 ukur-antrean.py | head -7

# Keluarannya:
#   Pekerjaan    Jenis                 Ke   Detik  Hasil
#   kerja-0001   kirim-laporan          1    0.01  terkirim dan di-ack
#   kerja-0002   gagal-dua-kali         1    0.01  gagal, webhook 500
#   kerja-0003   crash-setelah-kirim    1    0.01  terkirim, crash sebelum ack
#   kerja-0002   gagal-dua-kali         2    0.51  gagal, webhook 500
#   kerja-0003   crash-setelah-kirim    2    0.51  duplikat ditahan
#   kerja-0002   gagal-dua-kali         3    1.52  terkirim dan di-ack
\end{lstlisting}

\begin{table}[!htb]
  \centering
  \caption{Contoh hasil Latihan 3 yang sudah terisi}
  \label{tab:c13-lat3-contoh}
  \begin{tabularx}{\textwidth}{Xlll}
    \hline
    \textbf{Jeda dasar} & \textbf{Percobaan 2} & \textbf{Percobaan 3} & \textbf{Selisih} \\
    \hline
    0,5 detik & 0,51 & 1,52 & 0,5 lalu 1,0 \\
    \hline
  \end{tabularx}
\end{table}

\begin{table}[!htb]
  \centering
  \caption{Lembar isian hasil pengukuran Latihan 3}
  \label{tab:c13-lat3-hasil}
  \begin{tabularx}{\textwidth}{Xlll}
    \hline
    \textbf{Jeda dasar} & \textbf{Percobaan 2} & \textbf{Percobaan 3} & \textbf{Selisih} \\
    \hline
    0,5 detik & \dots & \dots & \dots \\
    \hline
    1,0 detik & \dots & \dots & \dots \\
    \hline
  \end{tabularx}
\end{table}

\ifajar
  \subsection*{Latihan 4: Duplikat dan Consumer Idempoten}
\else
  \subsection*{Prosedur 4: Duplikat dan Consumer Idempoten}
\fi

Latihan ini membuktikan bahwa pengiriman dapat terjadi dua kali, dan bahwa
akibat gandanya hanya dapat ditahan \textit{consumer}.

\begin{enumerate}
  \item Dari hasil Latihan 3, catat jumlah percobaan \textit{worker} dan
        jumlah pengiriman yang tiba di \textit{webhook}.
  \item Catat pekerjaan mana yang dikirim dua kali oleh \textit{worker}
        tetapi hanya sekali tiba di \textit{webhook}.
  \item Matikan pemeriksaan duplikat dengan memberi komentar pada panggilan
        \texttt{sudah\_diproses} di \berkas{worker.py}, jalankan ulang, lalu
        catat jumlah pengiriman yang tiba sekarang.
  \item Pulihkan kembali pemeriksaannya.
  \item Isi Tabel \ref{tab:c13-lat4-hasil}. Tabel \ref{tab:c13-lat4-contoh}
        memperlihatkan contoh yang sudah terisi.
  \item Jawab tiga pertanyaan berikut. Pertama, mengapa id pekerjaan harus
        tetap pada tiap percobaan? Kedua, mengapa pemeriksaan duplikat
        diletakkan sebelum pengiriman, bukan sesudahnya? Ketiga, pekerjaan
        apa saja yang akibatnya memang tidak berubah walaupun dikerjakan dua
        kali?
\end{enumerate}

\begin{lstlisting}[language=bash, caption={Membandingkan percobaan dengan pengiriman yang tiba, tersedia di \berkas{sources/chapter13/ukur-antrean.py}}, label={lst:c13-lat4-jalan}]
python3 ukur-antrean.py | tail -2

# Keluarannya:
#   Percobaan worker: 6, pengiriman tiba: 5
#   Pekerjaan tersisa: 0, ditandai diproses: 3
\end{lstlisting}

\begin{table}[!htb]
  \centering
  \caption{Contoh hasil Latihan 4 yang sudah terisi}
  \label{tab:c13-lat4-contoh}
  \begin{tabularx}{\textwidth}{Xlll}
    \hline
    \textbf{Keadaan} & \textbf{Percobaan} & \textbf{Tiba} & \textbf{Akibat ganda} \\
    \hline
    pemeriksaan duplikat hidup & 6 & 5 & tidak ada \\
    \hline
    \texttt{kerja-0003} saja & 2 & 1 & tidak ada \\
    \hline
  \end{tabularx}
\end{table}

\begin{table}[!htb]
  \centering
  \caption{Lembar isian hasil pengamatan Latihan 4}
  \label{tab:c13-lat4-hasil}
  \begin{tabularx}{\textwidth}{Xlll}
    \hline
    \textbf{Keadaan} & \textbf{Percobaan} & \textbf{Tiba} & \textbf{Akibat ganda} \\
    \hline
    pemeriksaan duplikat hidup & \dots & \dots & \dots \\
    \hline
    pemeriksaan duplikat dimatikan & \dots & \dots & \dots \\
    \hline
  \end{tabularx}
\end{table}
